% IMMUTABLE RED-TEAM REPORT::HASH-LENGTH::2026-10-07
\documentclass[11pt]{article}
\usepackage[T1]{fontenc}
\usepackage[margin=1in]{geometry}
\usepackage{amsmath,amssymb,amsthm,booktabs,array,hyperref}
\setlength{\emergencystretch}{3em}
\newtheorem{theorem}{Theorem}
\newtheorem{lemma}[theorem]{Lemma}
\newtheorem{corollary}[theorem]{Corollary}
\newtheorem{definition}[theorem]{Definition}
\newcommand{\SAT}{\operatorname{SAT}}
\newcommand{\E}{\mathrel{\equiv}}
\title{Hash Length Does Not Supply the Missing WRAPPER Contraction\\
       Exact Buckets, Equality Replay, Randomness, and COVER-SHRINK1}
\author{Independent red-team audit}
\date{7 October 2026}

\begin{document}
\maketitle

\begin{abstract}
This report audits the proposal that short hashes, adaptive rehashing, or an
external random bitstream might supply the missing contraction in the SAT
WRAPPER program.  Exact verification makes a hash an index only: the number
of safely reusable states remains the number of true semantic congruence
classes.  Omitting verification instead requires the hash itself to be a
sound semantic quotient, which is the missing merge theorem.  We derive exact
bucket and equality-replay bounds, distinguish deterministic worst-case from
random-oracle expectation, give a restriction-cover countermodel, and apply
the results to the concrete COVER-SHRINK1 target.  No choice of hash length
alone establishes that target or improves the proved joint time-space
exponent.  P versus NP remains open.
\end{abstract}

\section{Status and the concrete target}

The sufficient construction target under audit is the following.  Assuming
$P=NP$, an exact source-matched procedure on every explicit CNF of encoded
length $M$ should output at most
\[
 q\le M^{1/3+o(1)}
\]
ordinary-CNF children, each of encoded length at most
$M^{2/3+o(1)}$, using total local, generation, and combination work
$M^{1+o(1)}$ and peak workspace $M^{2/3+o(1)}$.  The construction must be
recursively closed and may not invoke the assumed SAT algorithm on an
$M^{1-o(1)}$-length query.  We call this target \textsc{Cover-Shrink1}.

The downstream composition arithmetic is not at issue here.  The question is
whether hashing supplies the missing all-input bound on the number of exact
children or makes their discovery cheap.  The answer is negative unless an
additional semantic quotient theorem is supplied.

\section{Exact hash-indexed reuse}

Let $X_1,\ldots,X_Q$ be states generated at one scale.  Let $\E$ be an
equivalence relation strong enough to justify every proposed future reuse.
For a recursive WRAPPER this may be normalized residual equality,
isomorphism, or another continuation congruence; equality of the present SAT
answer need not suffice.  Let $R$ be the number of $\E$-classes represented
among the $Q$ states.  Let
\[
 H:\mathcal X\longrightarrow\{0,1\}^h,
 \qquad B=2^h.
\]

\begin{lemma}[HASH-INDEX-NO-SHRINK]
Suppose a state is reused only after an exact verifier accepts it as
equivalent to a retained representative.  If the verifier decides $\E$
completely, complete deduplication retains exactly $R$ representatives.  If
the verifier is merely sound, it may retain more than $R$.  In either case
the hash length gives no upper bound on $R$.
\end{lemma}

\begin{proof}
Hash equality only selects candidate representatives.  Exact verification
never identifies two distinct $\E$-classes.  A complete verifier eventually
identifies every pair in one class, leaving one representative per class.  A
sound incomplete verifier may fail to identify members of the same class,
leaving additional representatives.  Changing the candidate index does not
change the underlying quotient.
\end{proof}

\begin{lemma}[Sound merge-on-hash bound]
If equal hashes are merged without exact verification, soundness requires
\[
 H(x)=H(y)\quad\Longrightarrow\quad x\E y.
 \tag{1}
\]
If $R$ equivalence classes occur, then
\[
 2^h\ge R,
 \qquad
 h\ge\lceil\log_2R\rceil.
 \tag{2}
\]
\end{lemma}

\begin{proof}
Two states from distinct classes cannot share a hash under (1).  Selecting
one representative from each class therefore gives $R$ distinct hash values.
\end{proof}

For current SAT-answer equivalence, one bit is sufficient in principle:
$H(F)=\SAT(F)$.  Constructing that hash is exactly the semantic decision
problem, so its full acquisition cost must be charged.  For future
restriction or continuation equivalence there can be exponentially many
classes.  For example, the unit-clause formulas selecting different unique
assignments are all presently satisfiable but are separated by later
restrictions.  Hash length and hash acquisition are different resources.

\section{Bucket and equality-replay bounds}

Consider the chained-bucket implementation in which a new state is compared
with earlier representatives having the same hash.  Suppose the $Q$ states
are pairwise inequivalent and bucket $z$ contains $m_z$ states.

\begin{lemma}[BUCKET-REPLAY]
The number of rejected exact comparisons is
\[
 P(H)=\sum_{z=1}^{B}\binom{m_z}{2}.
 \tag{3}
\]
It satisfies
\[
 P(H)\ge
 \max\left\{0,\frac{Q^2/B-Q}{2}\right\},
 \tag{4}
\]
and the worst case is $P(H)=\binom Q2$.
\end{lemma}

\begin{proof}
Within a bucket, the $j$th inequivalent insertion is rejected against its
$j-1$ predecessors.  Summing gives (3).  Since
$\sum_zm_z^2\ge Q^2/B$, equation (4) follows from
$P=(\sum_zm_z^2-Q)/2$.  Putting every state in one bucket attains the stated
worst case.
\end{proof}

If one exact comparison costs $E(M)$, equality-replay work is
$P(H)E(M)$.  A cache holding at most $K$ representatives can limit local
comparison work to $O(QK E)$, but then it does not perform global
deduplication.  A cyclic stream of $K+1$ different states defeats an LRU
cache.  Cache misses preserve correctness by recomputation; they do not
establish a smaller exact quotient.

Storing all representatives costs $\Omega(QD)$ bits when each retained
descriptor has $D$ bits, in addition to $Qh$ tag bits and table overhead.
Short source references can replace full representatives only if every
regeneration and exact replay is charged.

\section{Fixed public hashes and adversarial fibers}

\begin{lemma}[PUBLIC-FIBER]
For a fixed map $H:\{0,1\}^L\to\{0,1\}^h$, some bucket contains at least
$2^{L-h}$ strings.  Therefore a $Q$-element all-collision set exists whenever
\[
 L-h\ge\log_2Q.
 \tag{5}
\]
To exclude every $Q$-fold collision on the full domain requires
$h>L-\log_2Q$; injectivity requires $h\ge L$.
\end{lemma}

\begin{proof}
This is the pigeonhole principle applied to $2^L$ inputs and $2^h$ buckets.
\end{proof}

The same statement for a restricted valid-state domain $\mathcal D$ replaces
$2^L$ by $|\mathcal D|$.  There are exponentially many distinct normalized
CNF encodings, so polynomially many hostile states remain available for a
short fixed digest.

\begin{lemma}[Long-prefix replay countermodel]
For any integer
\[
 p\le L-h-\log_2Q,
\]
there is a set of $Q$ equal-hash strings sharing a prefix of length $p$.
Consequently a streaming equality test can be forced to inspect
$\Omega(p)$ bits per rejected candidate.
\end{lemma}

\begin{proof}
Fix any prefix of length $p$.  Its extension set contains $2^{L-p}$ strings.
One of the $2^h$ hash buckets contains at least $2^{L-p-h}\ge Q$ of them.
They all share the fixed prefix.  Choose suffixes whose first difference
appears only after that prefix.
\end{proof}

For $h+\log Q=o(L)$ this gives $\Omega(L)$ equality replay.  CNF states can
realize the same pattern with a long shared clause prefix and distinguishing
suffix clauses.  This is a countermodel for chained exact equality replay,
not a lower bound against a secondary exact dictionary or an independently
proved canonical representation.

A fixed cryptographic digest does not alter the theorem.  A 64-, 128-, or
256-bit output is $O(1)$ asymptotically.  Computational collision resistance
is not mathematical collision freedom and cannot support an exact all-input
complexity proof.  Under the contradiction hypothesis $P=NP$, relying on the
hardness of ordinary polynomially checkable collision search is especially
inappropriate.

\section{Random-oracle expectation}

For independent uniform $h$-bit hashes, pair indicators give
\[
 \mathbb E[P]=\frac{\binom Q2}{2^h}.
 \tag{6}
\]
The no-collision probability obeys
\[
 \Pr[\text{no collision}]
 =\prod_{i=0}^{Q-1}\left(1-\frac{i}{2^h}\right)
 \le
 \exp\left(-\frac{Q(Q-1)}{2^{h+1}}\right).
 \tag{7}
\]
A union-bound guarantee of collision probability at most $\delta$ requires
approximately
\[
 h\ge2\log_2Q+\log_2(1/\delta)-O(1).
 \tag{8}
\]

Write
\[
 Q=M^{u+o(1)},\qquad E(M)=M^{v+o(1)},
 \qquad h=\kappa\log_2M+O(1).
\]
Reading or hashing all states already costs $M^{u+v+o(1)}$ when each state
has $M^{v+o(1)}$ relevant bits.  Equation (6) gives expected replay exponent
$2u-\kappa+v$.  Hence
\[
 C_{\rm expected}\ge
 \max\{u+v,\,2u-\kappa+v\}.
 \tag{9}
\]

At the COVER-SHRINK1 scale $u=1/3$ and $v=2/3$,
\[
 C_{\rm expected}\ge\max\{1,4/3-\kappa\}.
 \tag{10}
\]
Thus $h=(1/3)\log_2M$ is enough for linear expected filtering work in the
random-oracle model, and $h=(2/3)\log_2M$ makes the expected number of false
collisions constant.  Neither is a deterministic worst-case conclusion.  A
fixed public hash still admits the all-collision fiber, giving
$M^{4/3-o(1)}$ chained replay when equality costs $M^{2/3-o(1)}$.

\section{Adaptive rehashing and set-dependent perfect hashing}

\begin{lemma}[Public rehash sequence]
For a fixed sequence $H_1,\ldots,H_t$ of $h$-bit hashes, the combined
signature has $th$ bits.  Among $R$ equivalence classes, one combined-signature
fiber contains at least
\[
 \left\lceil\frac{R}{2^{th}}\right\rceil
 \tag{11}
\]
classes.  Excluding a common fiber of $Q$ classes requires
\[
 th>\log_2(R/Q).
 \tag{12}
\]
\end{lemma}

\begin{proof}
Apply the pigeonhole principle to the vector map
$(H_1,\ldots,H_t)$, whose range has size at most $2^{th}$.
\end{proof}

Repeated public rehashing therefore accumulates the information of one longer
fingerprint.  Random rehashing plus exact verification is a Las Vegas
dictionary and gives expected, randomized bounds.  Deterministic seed
enumeration must charge every failed pass and its collision certification.

An input-dependent minimal perfect hash can avoid collisions on the actual
$Q$ keys with short labels.  This does not contradict the public-fiber lemma:
the selected map depends on the set.  It indexes $Q$ distinct keys but does
not merge semantic classes.  Its construction must inspect the keys, and its
table, representatives, or regeneration interface must be charged.  The
number of exact children remains $R$.

\section{External random bitstreams}

An external IID stream has four distinct interpretations.
\begin{enumerate}
\item Merge-on-random-hash without verification has nonzero error and is not
an exact SAT decider.
\item Las Vegas hashing with exact verification is correct but yields only
expected or high-probability randomized resource bounds.
\item A favorable stream selected separately for each input or input length
is nonuniform advice.
\item One globally fixed infinite random stream is an oracle or infinite
advice unless it is computable.  If computable, it is a deterministic public
schedule and is again subject to worst-case fibers.
\end{enumerate}

If correctness holds for every stream, a deterministic endpoint application
also needs the resource bound for every stream, including an all-zero or
otherwise hostile stream.  If performance holds only for almost every stream,
the result does not meet the deterministic Williams endpoint.  Even a
qualitative derandomization consequence of $P=NP$ would not automatically
preserve the fine-grained time and space exponents required here.

\section{A summary-only restriction-cover obstruction}

The following countermodel applies to hash-selected restrictions or
backdoors.  It deliberately does not claim a lower bound against arbitrary
source-dependent semantic CNF transformations.

\begin{lemma}[SUMMARY-COVER]
For each $x\in\{0,1\}^n$, let $F_x$ be the unit-clause CNF whose unique
satisfying assignment is $x$.  Suppose a constructor sees only an $h$-bit
summary $H(F_x)$ and, for each summary value, outputs at most $q$
restrictions, each leaving at most $\ell$ variables free.  Exact coverage of
every $F_x$ requires
\[
 \boxed{h+\log_2q+\ell\ge n.}
 \tag{13}
\]
\end{lemma}

\begin{proof}
Some summary bucket contains at least $2^{n-h}$ of the $2^n$ distinct unique
witnesses.  Every formula in that bucket receives the same restriction
family.  Each restriction covers at most $2^\ell$ assignments, so the family
covers at most $q2^\ell$.  Exact coverage forces
$q2^\ell\ge2^{n-h}$, which is equivalent to (13).
\end{proof}

Under a standard encoding, these formulas have $M=\Theta(n\log n)$ bits.
The choices
\[
 h=O(\log M),\qquad q=M^{1/3+o(1)},
 \qquad \ell\le M^{2/3+o(1)}
\]
make the left side of (13) $o(n)$, so they cannot cover this family.  Avoiding
the obstruction within the summary-only architecture requires
\[
 h\ge n-o(n)=M^{1-o(1)},
\]
which is incompatible with retaining the summary inside
$M^{2/3+o(1)}$ workspace.  A constructor that continues to read and exploit
the full source escapes the lemma, but that additional source-dependent
acquisition is exactly what COVER-SHRINK1 must bound.

\section{Hash-length regimes at the target parameters}

Let $Q=M^{1/3+o(1)}$ and let exact state replay have length or cost
$L=M^{2/3+o(1)}$.  The conclusions are summarized below.

\begin{center}
\small
\begin{tabular}{p{0.18\linewidth}p{0.33\linewidth}p{0.38\linewidth}}
\toprule
Hash choice & What it can establish & Why it does not establish COVER-SHRINK1 \\
\midrule
$h=O(1)$ & Constant-size bucket index & Chained exact replay can cost
$Q^2L=M^{4/3-o(1)}$; a bounded cache avoids that cost only by forfeiting
global deduplication. \\
\addlinespace
$h=\kappa\log M$ & Expected replay exponent
$\max\{1,4/3-\kappa\}$ under uniform random hashing & Fixed public hashes
retain all-collision worst cases.  Exact verification leaves $R$ classes;
random expectation does not meet the deterministic endpoint. \\
\addlinespace
$h=o(L)$ & More buckets and smaller expected collision count & A fixed map
still has $Q$-fold fibers whenever $L-h\ge\log Q$. \\
\addlinespace
$h>L-\log Q$ & Excludes $Q$-fold fibers on the full $L$-bit domain & The tag
is essentially as long as the state.  Retaining $Q$ tags costs
$M^{1-o(1)}$ bits, and injectivity produces no semantic merge. \\
\addlinespace
Set-dependent perfect hash & Collision-free indexing of the actual set &
Construction and table storage must be charged; $Q$ distinct semantic keys
remain $Q$ children. \\
\bottomrule
\end{tabular}
\end{center}

If one representative is declared per hash bucket with
$h\le(1/3+o(1))\log M$, either exact verification separates false collisions,
leaving $R$ children, or equal hashes are merged and condition (1) must be
proved.  The latter is a source-checkable semantic quotient with charged
construction cost---the missing theorem itself.

Hashes can be useful infrastructure after an independent theorem proves
$R\le M^{1/3+o(1)}$ and supplies cheap exact canonicalization or verification.
They do not prove that class-count bound, reduce child encoding length, or
remove current-scale acquisition.

\section{Reproducible finite arithmetic audit}

The companion script checks the bucket convexity inequality and the fixed
rehash common-signature bound, and simulates uniform random buckets for the
target powers.  It performed 1,872 bucket checks and 48 common-signature
checks with zero errors.  Eight deterministic random seeds were used for each
reported parameter row.  These are arithmetic sanity checks, not asymptotic
SAT evidence.

\begin{description}
\item[Script:] \texttt{hash\_length\_redteam\_2026-10-07.py}
\item[Results:] \texttt{hash\_length\_redteam\_results\_2026-10-07.json}
\item[Script SHA-256:]\leavevmode\\
\nolinkurl{a7db88107f9c99b548319bcc289cd82fbf9e1cf2c0c50452d6022892c5a18b75}
\item[Results SHA-256:]\leavevmode\\
\nolinkurl{0ac585e0c3599ff6cda9db4d2a36399623067a86b8b06ebe4ac38e5d83b28972}
\end{description}

\section{Final disposition}

The hash-length proposal has three legitimate roles: routing to possible
representatives, auditing provenance, and reducing ordinary dictionary
overhead once exact equivalence is already available.  None supplies the
load-bearing implication
\[
 P=NP
 \Longrightarrow
 \text{a recursively closed, source-matched COVER-SHRINK1 realization}.
\]
The remaining arrow is still a proof that the generated states have at most
$M^{1/3+o(1)}$ safely mergeable semantic classes, and that their discovery,
ordinary-CNF construction, exact combination, input access, and peak live
state fit the $M^{1+o(1)}$-time and $M^{2/3+o(1)}$-space ledger.  Hashes may
implement that theorem after it is proved; they do not imply it.

\begin{thebibliography}{9}
\bibitem{williams2010}
R. Ryan Williams,
\emph{Alternation-Trading Proofs, Linear Programming, and Lower Bounds},
STACS 2010, Theorem 3.9.
\url{https://drops.dagstuhl.de/entities/document/10.4230/LIPIcs.STACS.2010.2494}

\bibitem{williamsfull}
R. Ryan Williams,
\emph{Alternation-Trading Proofs, Linear Programming, and Lower Bounds},
full draft, Appendix B, Theorem B.1.
\url{https://www.cs.cmu.edu/~ryanw/automated-lbs.pdf}
\end{thebibliography}

\end{document}
